\chapter{多模态生成}
\label{chap:vision-multimodal-generation}

“让玩具狗绕花瓶走一圈，碰到花瓶时发出轻响”看似只有一句要求，交付时却要同时回答几个问题。玩具狗在转身后是否仍是原来的样子，花瓶有没有真的受到碰撞，声音是否发生在接触时刻？如果还要把玩具狗做成可以转动观察的三维资产，背面和侧面也必须共同解释同一个对象。单幅画面自然，只完成了其中一部分。

本章按交付成果讨论视频、配声、联合音画、空间内容及视觉条件响应。它们可以组成同一制作流程，但允许改变的变量不同。固定视频配声只能改变声音，联合生成可以同时调整动作和声轨，依据多幅观测拟合空间内容又有明确的图像证据。开始制作前，应记录给定条件、待生成内容和需要保持的关系，再选择具有相应入口的模型。

生成过程沿用\ref{sec:vision-models-media-generation}节的潜表示、条件连接和采样接口。下文用$s\in[0,1]$表示生成路径参数，用$t$或$\tau$表示媒体中的时间，避免把“第几次去噪”与“视频第几秒”混在一起。单幅图像创建与编辑见\ref{chap:vision-image-generation}章；本章增加跨时间、跨输出和跨视角的约束。

统一模型也要按实际输出入口使用。例如\ref{sec:vision-models-showo2}节的Show-o2可以在共享主干中形成文字与视觉输出，但视频能力取决于具体变体的训练范围，也没有因此获得音频输出。对于先解释参考图、再制作视频的请求，可以沿用其理解与视觉生成接口，声音仍需另一个支持配声的入口。共享主干和多模型组合都要接受同样的成果核验。

\section{视频生成与编辑}
\label{sec:vision-multigen-video}

对视频的操作可以改变整段画面、已有片段中的局部内容，或尚未给定的时间范围。三者共享生成基础，却不能共用一句含混的“根据视频生成视频”作为任务定义。原片段、首尾帧、区域掩码和文字说明分别提供什么信息，需要在进入模型前确定。

\subsection{从条件创建视频}
\label{sec:vision-multigen-create-video}

文字规定场景和事件，参考图像提供主体外观，首帧还规定故事开始时的状态。若请求包含精确轨迹，则需要模型接收随时间变化的位置或运动条件；仅支持文字的模型没有获得这条数值轨迹。相同的“镜头绕行”，也可能由相机移动、对象旋转或两者共同运动形成，不能只检查画面里出现了转动。

一个教学制作记录可以规定四秒、每秒24帧、固定画幅，以及“玩具狗从花瓶左侧走到右侧，花瓶保持不动”。按每帧占据一个播放时间间隔的约定，需要96帧，第$i$帧的起始时刻为$(i-1)/24$秒。若只规定首尾采样时刻相距四秒，则包含两端会得到97帧。实际编码器还可能要求帧数符合特定分组规则，制作记录应把采样范围、容器播放时长和模型输入长度分别写清。

\subsubsection{Lumiere的整段时空生成与重叠窗口超分辨率}
\label{sec:vision-multigen-lumiere}

逐帧独立生成时，即使每帧都有玩具狗，也可能在相邻帧改变耳朵形状。先生成稀疏关键帧再插入画面，还可能无法判断两帧之间究竟走了半圈还是一圈。Lumiere让一次网络前向处理覆盖整段视频，在压缩的时空表示中组织运动，再恢复原来的时间分辨率。这里依据其2024年v2说明\scite{vision-lumiere2024}

模型在预训练图像U-Net中加入时间模块，并在空间下采样、上采样附近增加时间分辨率变换。较细尺度使用分解的时空卷积，最粗尺度使用时间注意力。把时间长度压短以后，远隔帧的信息能够在较小表示上共同参与计算；随后逐层上采样，恢复完整帧率的低分辨率序列。

训练从带描述的视频构造含噪序列，保持原图像模型权重固定，更新新增时间模块。时间下采样会混合或选择帧，因此整个网络的初始化不能简单解释为“每帧恰好等于原图像模型”。原方法用接近最近邻下采样和帧复制上采样的初始化，为时间变换提供起点。

使用时从整段噪声视频出发，反复运行时空网络并执行采样更新。一次前向能看完整段，不表示一次前向已经完成采样。论文采用的短片段设置是80帧、每秒16帧，基础生成分辨率为$128\times128$；这说明该设计所覆盖的时间范围，不能由此推出任意长度的叙事都能一次处理。

低分辨率片段形成后，还要增加空间细节。高分辨率网络受显存限制，只能处理较短时间窗口。若各窗口独立完成再拼接，边界两侧可能出现不同的毛色或光照。Lumiere在每一步采样中处理重叠窗口，再对重复覆盖的位置融合预测，使下一步各窗口从同一段全局状态重新取样。

设窗口$i$选择全局视频中的若干连续帧，选择算子为$\bm P_i$；把当前一步的窗口预测展平成向量$\bm y_i$。希望新的全局状态$\bm z$在每个窗口内都接近相应预测，可以求解
\begin{equation}
  \min_{\bm z}\sum_i\lVert\bm P_i\bm z-\bm y_i\rVert_2^2,
  \qquad
  \widehat{\bm z}
  =\left(\sum_i\bm P_i^\top\bm P_i\right)^{-1}
     \sum_i\bm P_i^\top\bm y_i\text{。}
  \label{eq:vision-window-consensus}
\end{equation}
这里假设所有全局位置至少被一个窗口覆盖，且$\bm P_i$只负责取出对应坐标。对$\bm z$求导并令梯度为零，就得到右侧的正规方程。矩阵$\sum_i\bm P_i^\top\bm P_i$是对角矩阵，每个对角元素等于该位置被覆盖的次数，因此解就是对重叠预测逐位置求平均。

例如八帧教学序列采用窗口1至5和4至8，第4、5帧被预测两次。若第4帧某个颜色分量的两次预测为0.4和0.6，合并值为0.5；只被覆盖一次的分量直接保留。图\ref{fig:vision-window-consensus}展示了覆盖计数的来源。这个平均发生在每步采样的状态协调中，不能用“最终把两段成片淡入淡出”代替。

\begin{figure}[htbp]
  \centering
  \begin{tikzpicture}[x=1cm,y=1cm,font=\figurefont\small,
  cell/.style={draw=BookInk,line width=0.7pt,minimum width=7mm,
    minimum height=7mm,inner sep=0pt},
  arr/.style={-{Stealth[length=2mm]},BookInk,line width=0.8pt}]
  \node[anchor=east] at (1.65,2.45) {帧位置};
  \node[anchor=east] at (1.65,1.6) {窗口一};
  \node[anchor=east] at (1.65,0.65) {窗口二};
  \node[anchor=east] at (1.65,-0.45) {覆盖次数};
  \foreach \i in {1,...,8} {
    \pgfmathsetmacro{\xx}{1.5+0.95*\i}
    \node at (\xx,2.45) {\i};
    \ifnum\i<6
      \node[cell,fill=FigureInput!15] at (\xx,1.6) {};
    \fi
    \ifnum\i>3
      \node[cell,fill=FigureOutput!15] at (\xx,0.65) {};
    \fi
    \ifnum\i=4
      \node at (\xx,-0.45) {2};
    \else
      \ifnum\i=5
        \node at (\xx,-0.45) {2};
      \else
        \node at (\xx,-0.45) {1};
      \fi
    \fi
  }
  \draw[BookAmber,dashed,line width=0.8pt] (4.85,0.15) rectangle (6.7,2.02);
  \draw[arr] (5.3,0.1) -- (5.3,-0.1);
  \node[anchor=west] at (1.65,-1.3) {同一位置：$(0.4+0.6)/2=0.5$};
\end{tikzpicture}

  \caption{重叠窗口的教学构造。两个五帧窗口在第4、5帧重叠，覆盖次数决定式\eqref{eq:vision-window-consensus}的逐位置平均。图中帧数用于手算，不是Lumiere的发布配置；原文空间超分辨率窗口在时间上重叠两帧。}
  \label{fig:vision-window-consensus}
\end{figure}
\FloatBarrier

让玩具狗绕花瓶走一圈时，可分别检查起终点是否连接、步态是否连贯、遮挡后外观是否保持。重叠融合能协调窗口边界，却不能证明模型理解了接触或三维形状。短片段内部连贯与多镜头之间的身份保持，仍是不同的验收对象。

\Needspace{125mm}
\subsubsection{Wan的文字与首帧条件生成}
\label{sec:vision-multigen-wan}

把整段视频压进潜空间，可以减少生成网络处理的位置数。Wan技术报告v2采用因果视频自编码器、文字编码器和基于流匹配的生成主干。文字生成模型从噪声潜表示形成视频，图像条件版本则额外接收首帧的局部表示及整体外观信息\scite{vision-wan2025}

视频自编码器在时间上压缩，在每个空间方向也压缩。对于符合其长度约定的$T$帧视频，潜时间长度为$1+(T-1)/4$，空间尺寸分别为原来的$1/8$，潜通道数为16。第一帧单独保留时间起点，后续帧按时间组压缩。这里的因果性属于自编码器，不表示生成主干只能从前向后逐帧预测。

以$T=17$、$256\times256$画面为尺寸示例，目标潜网格为$5\times32\times32\times16$。图像条件版本把首帧放在条件视频的第一个位置，其余帧填零，然后通过同一个视频编码器得到条件潜网格。零表示未提供内容，不能把它当成“未来画面应为黑色”的监督目标。

模型还接收已知位置掩码，约定1表示已提供帧，0表示待生成帧。掩码保留原帧级信息，再按时间压缩组重排到通道维，与条件潜网格对齐。目标含噪潜变量16通道、条件潜变量16通道和重排后的掩码4通道，合起来形成36通道的输入。此处是报告中图像条件路线的输入组成；具体首帧重复和分组边界应沿用相应自编码器实现，不能只把二值掩码当作普通灰度视频压缩。

首帧经CLIP图像编码器提取的特征还通过MLP形成整体条件，并由独立交叉注意力接入主干。条件潜网格保留空间位置，整体特征补充外观语义，两者不能互相替代。训练前期混合图像到视频、续接、首尾转换等掩码任务；报告中的任务微调阶段再加入图像编码器分支，加强具体入口的条件控制。

目标视频编码为干净潜变量后，按\ref{sec:vision-models-flow-matching}节构造中间状态和速度目标。训练更新生成网络，使速度依赖文字、条件潜网格和掩码。推理从目标噪声潜网格出发，在相同条件下求解生成轨迹，再由视频自编码器恢复帧序列。编码器可以分块解码较长序列，并不自动意味着生成网络已经学会长时间状态保持。

固定同一首帧，分别要求玩具狗转头、向前走和镜头环绕，可以观察首帧条件留下了什么自由度。转头要补出原来看不到的侧面，行走要决定腿部运动和接触，镜头环绕还要补背景与背面。首帧只给了起点，这些未知内容不会因“图生视频”这个任务名称而变成已知。实际使用要对应具体检查点和条件入口，不能把报告中的多种任务并成每个权重都支持的功能。

\subsection{已有视频内容的修改}
\label{sec:vision-multigen-edit-video}

编辑把原片段本身作为条件，还要明确哪些变化是任务要求，哪些变化应算误改。把花瓶换成金属杯主要改变对象外观；让玩具狗跳上桌子则改变动作、遮挡和接触，原背景中某些区域也会随之显露。保留范围必须根据请求判断，不能一律要求整段视频的所有像素不变。

\Needspace{165mm}
\subsubsection{Movie Gen Edit的分阶段编辑监督}
\label{sec:vision-multigen-movie-edit}

只有“文字—视频”配对，模型学到的是描述什么画面，尚未学到根据原视频执行修改。Movie Gen Edit在视频生成主干中拼接原视频潜变量与含噪目标潜变量，并接入编辑指令和任务类型。其2024年v1通过分阶段构造训练关系，逐步缩小单帧编辑与真实视频编辑之间的差距\scite{vision-moviegen2024}

新增视频条件进入输入投影，任务类型具有学习到的嵌入，一部分接入文字上下文，另一部分接入生成时刻的调制。原视频和目标视频使用同一时间自编码器，但职责不同：原视频提供要修改的内容，目标潜变量在生成路径上逐步变化。新增连接采用零初始化，其余权重由预训练视频生成模型提供起点。

第一阶段把图像编辑配对视为单帧视频编辑，并与文字视频生成交替训练。图像编辑样本含输入图、编辑指令和目标图；纯视频生成样本则以黑色视频作为空输入条件。混合这两类任务，使模型既能利用编辑关系，又继续学习视频运动。单帧编辑训练的时间位置还随机选择，避免所有图像编辑都只出现在序列起点。

单帧训练缺少“原片段连续输入”的经验。第二阶段因此使用两类多帧构造。一类对原图和编辑后图逐步施加相同仿射变换，形成成对序列。由于两侧使用同样的变换，模型可以学习在相同运动下保留编辑差异，但这类平面移动并不能提供真实关节运动。另一类从真实视频定位并分割对象，将相应区域涂成指定颜色，把分割任务转成“把这个对象标成某颜色”的编辑配对，补充真实时间变化。

这些目标有一部分仍来自模型生成，可能携带失真。第三阶段使用编辑回译：把真实视频$V$按指令$r$编辑为$\widetilde V$，再根据前后描述形成逆向指令$r^{-}$，用$(\widetilde V,r^{-},V)$训练。目标回到真实视频，合成缺陷主要留在输入侧，模型不必继续把自己的缺陷当作目标学习。

例如真实片段里是红花瓶，前向请求为“改成金属杯”，合成输入是带杯子的片段，逆向请求为“把金属杯改回红花瓶”。训练目标仍是原始花瓶视频。这里的$r^{-}$是语义上构造并筛选的修改要求，不是保证可逆的数学算子。删除对象可能同时丢失背景细节，逆向指令也可能不准确，因此原方法仍需过滤训练配对。

推理时，原视频、指令及相应任务条件进入模型，目标潜变量从噪声逐步形成，再解码为编辑片段。检查金属杯的材质之外，还要检查杯体是否在遮挡后恢复为同一形状、玩具狗是否被连带改变。连续编辑时保存每轮输入与结果，才能发现哪一轮开始出现身份或背景漂移。Movie Gen的生成、编辑和配声分别使用相应路线，不能把模型族名称当成一个同时完成全部输出的网络。

\subsubsection{Lumiere的掩码条件视频修补}
\label{sec:vision-multigen-video-inpainting}

若需要直接指出修改范围，可以使用时空掩码。Lumiere的视频修补路线把含噪目标视频、保留可见区域的条件视频和二值掩码沿通道拼接。RGB目标3通道、RGB条件3通道、掩码1通道，因此输入层从3通道扩展到7通道，再微调生成模型。

为便于讨论，本节令$m_{t,i,j}=1$表示允许生成的位置，0表示要保留的位置。训练时从完整视频$\bm X_t$构造条件
$\bm C_t=(\bm1-\bm M_t)\odot\bm X_t$，
掩码沿颜色通道广播；完整视频的含噪版本仍作为待预测对象。模型利用未遮区域与文字学习补全，整段时空处理让各帧补出的杯子共享上下文。其他实现可能把掩码语义反过来，接入时须转换，不能把同一个0或1在不同接口中按相同含义解释。

推理时输入原视频、覆盖花瓶的逐帧掩码和替换要求，采样后检查边缘、遮挡及未指定区域。掩码若没有随花瓶移动，模型有时会在旧位置补出杯子，却在新位置留下原花瓶。区域跟踪与掩码传播可利用\ref{chap:vision-video-recognition}章的方法，但传播质量也属于编辑条件的一部分。

若交付明确要求掩码外像素完全相同，可以在制作阶段规定硬合成
\begin{equation}
  \bm X_t^{\mathrm{out}}
  =(\bm1-\bm M_t)\odot\bm X_t+
     \bm M_t\odot\widehat{\bm X}_t\text{。}
  \label{eq:vision-edit-composite}
\end{equation}
该式在掩码为0处代数上保留原像素；这是明确的输出规则，不是模型训练自动提供的保证。硬边界可能暴露接缝，改变光照或物体形状也可能需要更大的编辑范围。若使用软边缘，混合区域就不再逐像素相同，必须相应修改验收范围。

动态照片也可用同类输入组织：把一张静态图复制到各个时间位置，只开放局部区域变化。例如让玩具狗摇尾巴而背景固定，掩码应包含尾巴摆动需要扫过的空间。复制帧提供静态背景，没有给出尾巴的唯一真实运动。

\subsection{视频时间范围的延伸与补全}
\label{sec:vision-multigen-extension}

给定历史片段向后续接，只有已发生的状态；给定首尾两侧，中间过程还受到终点限制。创作允许新增动作与事件，而\ref{chap:vision-video-restoration}章的插帧侧重依据已采样观测估计中间时刻。两者可能使用相似掩码，评价证据却不相同。

\subsubsection{Wan的已知帧掩码与时间条件扩展}
\label{sec:vision-multigen-wan-extension}

Wan的掩码接口把已知画面放入指定时间位置，其余位置留空。掩码标明哪些内容是输入条件，再与含噪目标潜变量一起送入模型。开头一段已知对应续接，首尾已知对应两端之间的转换，零散已知帧则形成其他时间补全条件。共同接口的训练混合这些任务，具体使用仍需对应任务微调后的权重。

图\ref{fig:vision-time-masks}使用八个抽样位置展示条件差别。续接给出前三个状态，后面可以朝多种方向发展；首尾条件只固定两端，中间有多种合法运动；零散条件还限制途中必须经过的状态。图中的格子是教学时间位置，不是把八帧直接当作Wan自编码器的合法输入长度。

\begin{figure}[htbp]
  \centering
  \begin{tikzpicture}[x=1cm,y=1cm,font=\figurefont\small,
  cell/.style={draw=BookInk,line width=0.8pt,minimum width=7mm,
    minimum height=7mm,inner sep=0pt},
  known/.style={cell,fill=BookBlue,text=white}]
  \node[anchor=east] at (1.65,2.9) {时间位置};
  \node[anchor=east] at (1.65,2.1) {续接};
  \node[anchor=east] at (1.65,1.1) {首尾转换};
  \node[anchor=east] at (1.65,0.1) {零散条件};
  \foreach \i in {1,...,8} {
    \pgfmathsetmacro{\xx}{1.5+0.95*\i}
    \node at (\xx,2.9) {\i};
    \ifnum\i<4
      \node[known] at (\xx,2.1) {1};
    \else
      \node[cell] at (\xx,2.1) {0};
    \fi
    \ifnum\i=1
      \node[known] at (\xx,1.1) {1};
    \else
      \ifnum\i=8
        \node[known] at (\xx,1.1) {1};
      \else
        \node[cell] at (\xx,1.1) {0};
      \fi
    \fi
    \ifnum\i=2
      \node[known] at (\xx,0.1) {1};
    \else
      \ifnum\i=5
        \node[known] at (\xx,0.1) {1};
      \else
        \ifnum\i=8
          \node[known] at (\xx,0.1) {1};
        \else
          \node[cell] at (\xx,0.1) {0};
        \fi
      \fi
    \fi
  }
  \draw[-{Stealth[length=2mm]},BookInk,line width=0.8pt]
    (2.1,-0.7) -- (9.5,-0.7) node[below left] {时间递增};
\end{tikzpicture}

  \caption{时间条件的三种组织。实心格和数字1表示已提供内容，空心格和数字0表示待生成内容。位置按时间排序，编码与潜空间重排发生在此后；同一掩码形式不保证任意发布权重都支持相应任务。}
  \label{fig:vision-time-masks}
\end{figure}
\FloatBarrier

让玩具狗离开画面后继续前行，模型可以决定离开后的路线；给定它最终位于桌右侧的终帧，就要求生成过程连接两个状态。判断连接质量时，除了比较边界帧外观，还可以记录图像位置的有限差分。若边界前两次横向位移为2、2像素，新片段第一次位移为9像素，即使边界图像接得上，也出现了明显的速度变化。透视和相机运动会影响像素速度，所以此数值只用于固定相机的教学诊断。

反复续接时，新片段成为下一轮条件，误差可能进入后续每一段。花瓶位置每次只漂移少量像素，长序列里也可能明显偏离原场景。首尾相接、局部运动连贯与全程状态一致应分别检查。把两端之间的创作画面用于说明一个故事是合理的；把它当成实际未拍到过程的观测证据则超出了条件所能支持的结论。

\section{视觉条件下的声音生成}
\label{sec:vision-multigen-audio}

本节固定已有视频，待生成对象是声轨。模型可以在训练时同时使用视频、文字与音频，但只要推理输出只改变声音，就仍是视觉条件下的配声任务。它与同时创造画面和声音使用不同的可修改范围。

\subsection{视觉事件与声音要求的确定}
\label{sec:vision-multigen-audio-conditions}

画面中一次接触可以对应轻敲、摩擦或破碎，具体声音还取决于材料、力度和录音位置。玩具狗移动也可以配脚步或机械声。视频往往没有唯一决定这些声学属性，因此需要把可见事件与创作选择分开记录。

可以用事件记录$(\tau_i^{\mathrm{start}},\tau_i^{\mathrm{end}},e_i,r_i)$描述第$i$个声音要求，其中时间以秒计，$e_i$表示对应事件，$r_i$说明声源、内容及保持要求。某次敲击记录来自视频中的接触；画外鸟鸣记录来自创作说明。两者都能成为有效条件，但后者不是模型从像素直接测得的事实。

同一四秒教学片段中，假定两次接触发生于1.0秒和2.5秒，可规定每次接触配一个短促轻响，其他时间保留低音量环境声。这份记录仍不等于模型一定支持逐事件数值控制；如果入口只有文字和视频，它只是验收依据。只有模型实际接收时间标记、局部掩码或其他对应条件时，这些字段才直接进入生成计算。

\subsection{声音生成与时间对应}
\label{sec:vision-multigen-audio-alignment}

语义特征帮助决定“什么声源”，时间特征帮助决定“何时发声”。一段视频的平均特征可以保留花瓶和桌面，却难以说明两次敲击之间的停顿。生成器必须同时利用场景内容与局部变化，才能把声轨安排到画面事件上。

\subsubsection{MMAudio的联合声文训练与帧级同步条件}
\label{sec:vision-multigen-mmaudio}

MMAudio把视觉语义特征、可选文字和随时间变化的同步特征作为音频生成条件。其多模态联合训练结合音视频配对与音频文字配对，输出仍是音频潜变量。本节依据论文v2解释这一接口\scite{vision-mmaudio2025}

音频先转为梅尔频谱，即按听觉相关频带汇总的短时频谱，再由预训练自编码器压成潜序列。视觉和文字经CLIP提取特征。生成主干的前一组模块让文字、视觉与含噪音频共同参与注意力，后一组模块只更新音频表示，最后预测音频潜变量的速度。视觉与文字的中间特征可以在网络内变化，原视频帧仍保持固定。

三种序列不具有相同的时间分辨率。原论文默认16 kHz音频版本的潜序列为每秒31.25个位置，视觉语义特征为每秒8个位置。把二者相同的序号直接当作同一时刻会发生错位。模型按比例调整旋转位置编码，使视频中经过一秒所产生的位移，与音频潜序列中经过一秒具有对应尺度。

对于更细的同步信息，模型使用在音画偏移检测上预训练的Synchformer视觉编码器。它读取短时视频，提供与事件发生时间相关的特征。原文主要八秒和十秒设置中的有效同步特征率为每秒24个位置；附录的提取过程从每秒25帧输入形成重叠短片段，不能把24直接当作所有输入长度下不变的原始取样率。

同步特征经过局部卷积处理，并用最近邻插值对齐到音频潜位置，再加上由文字、视觉及生成时刻得到的全局条件。令对齐后的第$j$个条件为$\bm c_j$，该位置音频隐状态为$\bm h_j\in\mathbb R^d$，调制写成
\begin{equation}
  \widetilde{\bm h}_j
  =\bm\gamma(\bm c_j)\odot
     \operatorname{LayerNorm}(\bm h_j)
    +\bm\beta(\bm c_j)\text{。}
  \label{eq:vision-audio-sync-modulation}
\end{equation}
两个学习映射给出逐通道缩放与偏移。全局条件会在所有$j$处重复，同步条件则随$j$变化，因此第1秒附近可以得到与第2.5秒不同的调制。这里没有把同步编码器的偏移分类结果直接当作波形，它提供的是生成条件。

在默认16 kHz设置中，八秒音频包含128000个波形采样点，压缩后有250个潜位置；同一段视频的语义特征有64个位置，同步特征有192个位置。对齐不是要求这三个长度相等，而是让它们按共同时间关系进入交互。插值不会新增视频中未观察到的敲击事件，特征本身漏掉动作时，调制也不能凭空恢复证据。

训练把真实音频潜变量与噪声连接，回归条件速度，复用式\eqref{eq:vision-flow-loss}的目标。音频文字数据缺少视频时，视觉与同步入口使用学习到的空条件，使模型仍能学习声音分布及文字关系。推理固定视频和文字，从音频噪声积分到干净潜表示，再由自编码器恢复频谱、由声码器恢复波形。声码器负责把频谱变成可播放采样，不负责判断画面中的碰撞是否发生。

固定花瓶落桌片段，分别给出“轻轻敲击”和“陶瓷破碎”的文字，可以区分内容选择与时间对应。若画面没有破碎而声轨出现碎裂声，作品可能满足了一个夸张音效要求，却不符合写实碰撞要求。MMAudio的音效和环境声结果也不能自动证明长段对白可懂、说话人稳定或歌词准确，这些输出需要自己的训练范围与评价。

\subsubsection{Movie Gen Audio的上下文声轨延伸与修补}
\label{sec:vision-multigen-movie-audio}

声音跨片段还需保持背景、响度与事件连续。前面引用的Movie Gen报告在音频路线中使用掩码音频上下文，学习根据视频、文字及已知声音预测待补部分。没有上下文对应整段生成，已知一侧对应延伸，两侧已知对应局部修补。

该路线用单独训练的DAC-VAE将波形压成连续潜序列，保留内容的位置作为条件，其余位置填零。条件潜变量与含噪目标潜变量逐位置拼接。视频每帧的语义特征按最近邻方式对齐到音频潜时间尺度，经投影后与音频表示相加；文字通过交叉注意力提供声源和音乐等要求。训练仍回归流匹配速度，生成后用音频解码器恢复波形。

长声轨可以按片段逐段产生，把上一段末尾的音频作为下一段的已知上下文。若两个片段都包含同一次碰撞，就不能在最终结果里各保留一次，必须按共同时间位置合并重叠区域。原方法的逐段路线还在生成路径的每一步混合重叠潜状态，使当前片段不过度偏离上一段的对应部分。

报告另采用多窗口共同更新的路线：每次生成步把全局音频潜序列分成重叠窗口，各自预测，再按窗口权重归一化合并，之后进入下一步。它与\ref{sec:vision-multigen-lumiere}节的协调方式相近。用逐渐增减的窗口权重，可以减少从单窗口突然切到等权混合造成的突变；若某个位置的权重总和为零，则无法归一化，实际边界必须保证覆盖。

教学上，先为固定视频形成环境声，再只替换第二次敲击。条件声轨在第二次敲击附近留空，其余部分作为上下文；生成后检查新敲击是否出现、周围环境是否连续。即使潜空间有上下文条件，解码器仍可能影响邻近波形，所以要求严格保持的区间要在最终波形上核验。报告的原音频任务不包含语音或带人声的音乐，不能从器乐与音效生成推断对白能力。

\subsection{声轨的核验与调整}
\label{sec:vision-multigen-audio-check}

声音可以“听起来像敲击”，却早于画面半秒出现。对短促事件，可以先匹配声源与次数，再比较已匹配事件的起始时刻。设视频接触时刻为$\tau_i^{\mathrm v}$，音频起声时刻为$\tau_i^{\mathrm a}$，定义
\begin{equation}
  \delta_i=\tau_i^{\mathrm a}-\tau_i^{\mathrm v},
  \qquad
  e_{\mathrm{sync}}=\frac1m\sum_{i=1}^{m}|\delta_i|\text{。}
  \label{eq:vision-event-sync-error}
\end{equation}
$\delta_i>0$表示声音较晚。这里的$m$是已正确配对事件数，漏声、多声与声源错误另行计数，不能只对少量成功配对求平均而隐藏其余失败。

沿用1.0秒和2.5秒的接触，若起声时刻为1.12秒和2.48秒，偏移为120毫秒和$-20$毫秒，平均绝对偏移70毫秒。若教学验收预先规定每次不得超过80毫秒，则只有第二次通过。这是便于计算的任务阈值，不是普适的人类听觉阈值。图\ref{fig:vision-event-sync}把两种方向的偏移放到共同时间轴上。

\begin{figure}[htbp]
  \centering
  \begin{tikzpicture}[x=1cm,y=1cm,font=\figurefont\small,
  axis/.style={BookInk,line width=0.8pt},
  contact/.style={BookBlue,line width=1.2pt},
  sound/.style={BookAmber,line width=1.2pt}]
  \node[anchor=east] at (1.6,3.4) {接触};
  \node[anchor=east] at (1.6,2.6) {起声};
  \draw[axis] (2,3.4) -- (10,3.4);
  \draw[axis] (2,2.6) -- (10,2.6);
  \foreach \t in {0,...,4} {
    \pgfmathsetmacro{\xx}{2+2*\t}
    \draw[axis] (\xx,3.35) -- (\xx,3.45);
    \draw[axis] (\xx,2.55) -- (\xx,2.65);
    \node[below] at (\xx,2.45) {\t};
  }
  \node[anchor=west] at (7.45,1.78) {媒体时间（秒）};
  \foreach \xx in {4,7} {
    \draw[contact] (\xx,3.2) -- (\xx,3.65);
    \fill[BookBlue] (\xx,3.4) circle (1.5pt);
  }
  \foreach \xx in {4.24,6.96} {
    \draw[sound] (\xx,2.6) -- ++(-0.1,-0.16) -- ++(0.2,0) -- cycle;
  }
  \node at (4,3.95) {1.00};
  \node at (7,3.95) {2.50};
  \node[anchor=west] at (0.3,1.3) {第一次：声音较晚};
  \node[anchor=west] at (5.75,1.3) {第二次：声音较早};
  \draw[axis,-{Stealth[length=2mm]}] (0.6,0.1) -- (5,0.1);
  \draw[axis,-{Stealth[length=2mm]}] (6,0.1) -- (10.4,0.1);
  \draw[contact] (1.7,-0.08) -- (1.7,0.35);
  \draw[sound] (3.9,-0.08) -- (3.9,0.35);
  \draw[sound] (7.4,-0.08) -- (7.4,0.35);
  \draw[contact] (8.35,-0.08) -- (8.35,0.35);
  \draw[BookInk,-{Stealth[length=2mm]},line width=0.8pt]
    (1.7,0.58) -- (3.9,0.58);
  \draw[BookInk,-{Stealth[length=2mm]},line width=0.8pt]
    (8.35,0.58) -- (7.4,0.58);
  \node[above] at (2.8,0.62) {$+120$ ms};
  \node[above] at (7.88,0.62) {$-20$ ms};
  \node[below] at (1.7,-0.12) {1.00};
  \node[below] at (3.9,-0.12) {1.12};
  \node[below] at (7.25,-0.12) {2.48};
  \node[below] at (8.5,-0.12) {2.50};
  \node[anchor=west] at (0.3,-0.85) {局部放大分别使用自己的时间尺度，刻度单位为秒。};
\end{tikzpicture}

  \caption{两次碰撞的教学时间对齐。上排为可见接触，下排为音频起声，局部放大给出偏移方向与大小；不是模型实测波形。平均值可能掩盖第一次事件超出既定容差，漏声和多声还须单独记录。}
  \label{fig:vision-event-sync}
\end{figure}
\FloatBarrier

持续环境声没有唯一冲击起点，应用场景与前后连贯性更合适；对白还需要转写内容、说话人和口型关系。近景同步测量也要说明是否扣除音频设备延迟，远距离声源则可能具有合理的传播延迟。声画不同时出现并不总是错误，判断取决于所描述的物理或创作条件。

调整要针对失败来源。声源不合适可改文字，局部时间错误可重做相应窗口，长背景突变可改用上下文延伸或修补。每次保存改动范围和候选，重新检查与邻段的连接。修改、延伸与修补是不同条件下的选择，不是所有配声都必须经历的固定步骤。

\section{音视频的协同生成}
\label{sec:vision-multigen-joint}

当画面与声音都尚未固定，可以先决定一边再生成另一边，也可以让两边在同一生成过程中交换信息。两种路线都要落到共同事件和时间上，区别在于生成过程中哪一边还有调整空间。

\subsection{分阶段生成与条件传递}
\label{sec:vision-multigen-staged}

先制作玩具狗碰倒花瓶的视频，再按成片配声，可以复用前两节的方法。视频阶段交付的不只是帧文件，还应包括实际时长、帧率、事件时间、可见声源以及声音可自由设计的部分。若碰撞从原计划的1.0秒变成1.4秒，配声条件应采用成片中的时间。

制作过程中可把理想化的分布关系写成
\[
  p(V,A\mid c)=p(V\mid c)\,p(A\mid V,c)\text{，}
\]
其中$V$为视频，$A$为音频，$c$为创作条件。这是概率链式分解，本身没有损失一般性；实际误差来自两个模型的近似、条件传递以及第二阶段受已选视频约束。不能仅从分阶段形式推出它在数学上无法表示联合分布。

在固定成片上配声时，声音模型只能适应已有碰撞。若画面没有真正接触花瓶，补一声碰撞不会让动作成立。此时应回到视频阶段修改，再更新配声，或者选择允许共同调整的模型。不同模型间共享时间单位和事件记录，可以组成一致的制作流程，但不等于它们共享参数或共同训练。

\subsection{联合生成与相互约束}
\label{sec:vision-multigen-joint-mechanism}

联合生成把画面和声音都放在待采样状态中。文字应同时说明可见动作与声音要求，训练配对要来自同一时间范围，否则分支之间学到的联系可能是错误的。网络交互提供了表达共同事件的通道，其效果仍需检查最终音画。

\subsubsection{Ovi的双分支信息交换与共同时间采样}
\label{sec:vision-multigen-ovi}

Ovi的v1使用结构相配的两个生成主干，一个处理视频潜表示，另一个处理音频潜表示。每层通过双向交叉注意力交换信息，共享冻结T5编码器产生的文字条件。文字请求将视觉描述、可听事件及对白内容组织在同一个上下文中\scite{vision-ovi2025}

相同隐藏宽度不表示两个原始潜空间相同。视频有时间与空间位置，音频以时间序列组织，进入主干前各有相应映射。交叉注意力读取另一模态的隐藏特征。论文的音频和视频潜序列长度不同，使用按比例缩放的旋转位置编码建立时间对应，不能让第31个音频位置无条件对应第31个视频位置。

训练先以音频文字数据建立音频主干，学习语音和音效，再把它与预训练视频主干连接。联合阶段新建模态间注意力，更新注意力模块并冻结前馈网络；配对音视频提供两种干净目标，各自抽取独立噪声，但使用同一个生成时刻$s$。两个速度目标分别是自身的数据端点减噪声端点，共同损失按权重相加。

推理时，把联合状态视为一个较长的向量，其演化可写为
\begin{equation}
 \frac{\mathrm d}{\mathrm ds}
 \begin{pmatrix}\bm z^{\mathrm v}\\ \bm z^{\mathrm a}\end{pmatrix}
 =
 \begin{pmatrix}
 \bm v^{\mathrm v}(\bm z^{\mathrm v},\bm z^{\mathrm a},s,c)\\
 \bm v^{\mathrm a}(\bm z^{\mathrm a},\bm z^{\mathrm v},s,c)
 \end{pmatrix}\text{。}
 \label{eq:vision-joint-flow}
\end{equation}
每边速度都依赖两边当前状态。采用欧拉法说明时，两个下一步都应由同一个当前状态计算，再共同更新；若先把视频推进一步，然后用新视频与旧音频计算音频速度，就已换成另一种分步更新。原论文使用共同日程和ODE求解器，并非只做一次网络前向。

教学上取单维潜状态$z^{\mathrm v}=0.2$、$z^{\mathrm a}=-0.1$，假设当前预测速度分别为0.6和0.4，步长0.25。一次共同欧拉更新得到0.35和0。这个计算只说明数值更新的同步，没有证明媒体里的敲击已同步。生成时刻$s$走到1以后，还需分别解码视频和音频，按媒体时间组合播放。

请求中可以同时规定“玩具狗接触花瓶”“短促金属碰撞声”和“安静房间环境”。结果应检查声源、动作及起止是否属于同一事件。说话场景还要检查文字、音素和嘴部运动，不能用敲击同步代替口型验证。Ovi原版聚焦约五秒片段；将多段分别成功的片段拼成一分钟故事，仍需另外解决跨镜头身份、场景与叙事状态。

\subsubsection{LTX-2的模态条件引导与联合结果形成}
\label{sec:vision-multigen-ltx}

LTX-2的两种媒体表示及双向连接已在\ref{sec:vision-models-ltx2}节说明。其2026年v1在推理时还允许分别增强文字条件和另一模态带来的影响。使用差异发生在每一步的预测组合中，不是生成完成后才为声音与画面打分。

固定当前模态的潜状态及生成时刻，记$\bm f(c,m)$为同时给定文字$c$和另一模态信息$m$时的预测，$\bm f(\varnothing,m)$为去掉文字的预测，$\bm f(c,\varnothing)$为去掉另一模态信息的预测。报告的引导组合可写为
\begin{equation}
 \widehat{\bm f}
 =\bm f(c,m)
  +\lambda_c[\bm f(c,m)-\bm f(\varnothing,m)]
  +\lambda_m[\bm f(c,m)-\bm f(c,\varnothing)]\text{。}
 \label{eq:vision-modality-guidance}
\end{equation}
两项差分别刻画加入文字与加入另一模态后当前预测如何变化。这里把完整条件预测作为基准，故$\lambda_c=\lambda_m=0$仍保留两种条件，不能误读为无条件生成。去条件分支也应使用方法规定的缺省入口，不是任意填一幅黑图就等价。

假设某个潜分量三次预测依次为0.6、0.2、0.5，取$\lambda_c=2$、$\lambda_m=1$，组合值为$0.6+2(0.6-0.2)+(0.6-0.5)=1.5$。它是潜空间预测，可以超出三次预测的范围，不是必须在0至1之间的像素。提高引导会改变采样轨迹，同时增加去条件分支的计算，是否改善当前请求需比较完整结果。

同一“玩具狗靠近花瓶并发出机械声”请求，可以固定初始噪声、长度与采样步数，改变$\lambda_m$，查看声音是否更贴近动作、画面是否出现额外变化。随后再在多组噪声下重复，避免单个候选主导判断。若文字与首帧冲突，增强两个条件也不自动解决冲突，需要先明确优先保留哪一个。

高分辨率制作可以先生成基础潜变量，再做空间潜上采样，并对重叠时空区块细化后合并。两路潜变量最终分别恢复为视频帧与音频频谱，频谱经声码器形成波形。容器层仍要写入正确时长、采样率与时间戳；如果导出时错误改变帧率，模型内部的时间联系也会被破坏。分块降低资源压力，不提供无限长叙事的一致性保证。

\section{多视角与三维内容生成}
\label{sec:vision-multigen-3d}

空间内容需要从不同方向解释同一个对象。无参考时，外观主要由文字和生成先验决定；单幅参考约束可见面；多幅或连续观测进一步约束空间与运动。三者都可能输出渲染图，但图像展示、可编辑网格和物理模拟模型有不同的验收要求。

\subsection{无参考观测下的空间内容创建}
\label{sec:vision-multigen-no-reference}

文字“陶瓷玩具狗”没有给出目标对象的真实背面。生成器可以创作一个符合要求的对象，但各方向不能各自长出一个正面。让所有相机都通过同一个三维表示形成图像，是把二维生成先验联系到空间内容的一种方法。

\Needspace{50mm}
\subsubsection{DreamFusion的二维先验引导三维优化}
\label{sec:vision-multigen-dreamfusion}

DreamFusion为每条文字请求从随机初始化的三维辐射场出发，反复渲染随机相机视图，使用冻结的二维扩散模型提供修改方向。它的\term{分数蒸馏采样}（score distillation sampling，SDS）把二维先验偏好的变化传回三维参数，使不同视角的图像逐步接近文字要求\scite{vision-dreamfusion2023}

辐射场与体渲染见\ref{sec:vision-generation-nerf}节。这里设场景参数为$\bm\theta$，相机和光照条件为$\pi$，渲染结果展平为$\bm x=g_{\bm\theta}(\pi)\in\mathbb R^d$。噪声水平用$\ell$表示，加入噪声得到$\bm z_\ell=\alpha_\ell\bm x+\sigma_\ell\bm\epsilon$，其中$\bm\epsilon\sim\mathcal N(\bm0,\bm I)$。冻结模型在文字$c$下预测噪声$\widehat{\bm\epsilon}_\phi$。

SDS给出的参数梯度估计为
\begin{equation}
 \bm g_{\mathrm{SDS}}
 =w(\ell)\bm J_g^\top
     [\widehat{\bm\epsilon}_\phi(\bm z_\ell,\ell,c)
       -\bm\epsilon],
 \qquad
 \bm J_g=\frac{\partial g_{\bm\theta}(\pi)}
                 {\partial\bm\theta}\text{。}
 \label{eq:vision-sds-gradient}
\end{equation}
若参数数目为$p$，$\bm J_g\in\mathbb R^{d\times p}$，噪声差为$d$维，结果是$p$维参数梯度。优化器沿其负方向更新场景。二维模型参数$\phi$保持固定，计算中也不对噪声预测网络的输入Jacobian反向传播。

后一限制不是普通噪声平方误差求导后的自然结果。若直接最小化
$\|\widehat{\bm\epsilon}_\phi(\bm z_\ell,\ell,c)-\bm\epsilon\|^2$，
链式法则还会出现噪声网络对$\bm z_\ell$的导数。SDS通过概率分布之间的目标获得另一种方向，可以由下面的简化推导看清。

固定$\pi$和$\ell$，令$q_{\bm\theta}$为均值$\alpha_\ell g_{\bm\theta}(\pi)$、方差$\sigma_\ell^2\bm I$的高斯分布，$p_\phi(\bm z_\ell\mid c)$为扩散先验在该噪声水平的条件分布。因为$q_{\bm\theta}$的方差不变，熵不随均值改变。利用重参数化，
\[
 \nabla_{\bm\theta}\operatorname{KL}(q_{\bm\theta}\|p_\phi)
 =-\alpha_\ell\bm J_g^\top
     \mathbb E_{\bm\epsilon}
       [\nabla_{\bm z_\ell}\log p_\phi(\bm z_\ell\mid c)]\text{。}
\]
用噪声网络提供的分数近似
$\nabla_{\bm z_\ell}\log p_\phi\approx
-\widehat{\bm\epsilon}_\phi/\sigma_\ell$，
再乘以权重$w(\ell)\sigma_\ell/\alpha_\ell$，得到
$w(\ell)\bm J_g^\top\mathbb E[\widehat{\bm\epsilon}_\phi]$。
固定相机和噪声水平时，$\bm J_g$不依赖抽样噪声，而$\mathbb E[\bm\epsilon]=\bm0$，因此可减去$\bm\epsilon$作为零均值控制变量，形成式\eqref{eq:vision-sds-gradient}的估计。

这一推导要求所用噪声水平的$\alpha_\ell,\sigma_\ell$非零，并把学习到的分数当作分布分数的近似。实际引导和网络误差会影响方向，不能因此声称每步都精确降低真实三维分布的KL散度。实现可对$w(\ell)(\widehat{\bm\epsilon}_\phi-\bm\epsilon)$停止梯度，再让它与渲染图像作内积求导；直接对噪声差的平方求导会得到不同计算。

用一个局部线性教学渲染器检查方向。设$\bm\theta=(a,b)^\top$，两个像素分量为$g_{\bm\theta}=(a,a+b)^\top$，则
$\bm J_g=\left(\begin{smallmatrix}1&0\\1&1\end{smallmatrix}\right)$。
若权重为1、噪声差为$(0.2,-0.1)^\top$，参数梯度为$(0.1,-0.1)^\top$。从$(0.5,0.2)^\top$以步长0.1更新，得到$(0.49,0.21)^\top$，渲染变为$(0.49,0.70)^\top$。第二像素同时受$a,b$影响，两者变化抵消，说明像素方向不能不经渲染Jacobian直接当作参数方向。图\ref{fig:vision-sds-path}标出了梯度经过和停止的位置。

\begin{figure*}[htbp]
  \centering
  \begin{tikzpicture}[x=1cm,y=1cm,font=\figurefont\small,
  block/.style={draw=FigureModel,fill=FigureModel!10,line width=0.8pt,
    minimum height=10mm,minimum width=23mm,align=center},
  input/.style={block,draw=FigureInput,fill=FigureInput!10},
  arr/.style={-{Stealth[length=2mm]},BookInk,line width=0.8pt,
    rounded corners=1.2mm},
  grad/.style={-{Stealth[length=2mm]},FigureLoss,dashed,line width=1pt,
    rounded corners=1.2mm}]
  \node[input] (theta) at (1.2,2) {三维参数\\$\bm\theta$};
  \node[block] (render) at (4.5,2) {渲染器\\$g_{\bm\theta}(\pi)$};
  \node[block] (noise) at (7.8,2) {加噪\\$\alpha_\ell\bm x+\sigma_\ell\bm\epsilon$};
  \node[block] (prior) at (11.1,2) {冻结二维先验\\$\widehat{\bm\epsilon}_\phi$};
  \node[block] (res) at (14.6,2) {噪声差\\$\widehat{\bm\epsilon}_\phi-\bm\epsilon$};
  \draw[arr] (theta) -- (render);
  \draw[arr] (render) -- node[above] {$\bm x$} (noise);
  \draw[arr] (noise) -- (prior);
  \draw[arr] (prior) -- (res);
  \node (camera) at (4.5,3.5) {相机与光照$\pi$};
  \node (eps) at (7.8,3.5) {噪声$\bm\epsilon$};
  \node (text) at (11.1,3.5) {文字$c$};
  \draw[arr] (camera) -- (render);
  \draw[arr] (eps) -- (noise);
  \draw[arr] (eps.north) -- (7.8,4.25) -- (14.6,4.25) -- (res.north);
  \draw[arr] (text) -- (prior);
  \node[block,draw=FigureLoss,fill=FigureLoss!8,minimum width=35mm]
    (jac) at (4.5,-0.15) {通过渲染器求导\\$w(\ell)\bm J_g^\top(\widehat{\bm\epsilon}_\phi-\bm\epsilon)$};
  \draw[grad] (res.south) -- (14.6,-0.15) -- node[above] {噪声差停止梯度} (jac.east);
  \draw[grad] (jac.west) -- (1.2,-0.15) -- (theta.south);
  \node[anchor=west] at (7.1,-0.8) {不反传噪声网络的输入Jacobian};
\end{tikzpicture}

  \caption{SDS把冻结二维先验的噪声差传回三维参数。实线为渲染、加噪与预测，虚线为通过渲染器的参数梯度。噪声差在求参数梯度时视为常量，不反传冻结噪声网络的输入Jacobian；每次重新抽取相机，使多个方向约束同一个场景。}
  \label{fig:vision-sds-path}
\end{figure*}
\FloatBarrier

原方法每次抽样相机和光照，渲染密度与反照率，并结合表面朝向等正则项优化场景。反照率表示材料本身的颜色倾向，光照与法线参与当前图像形成；这有助于把部分外观变化与几何区分。模型还按相机方向调整文字中的视角描述，减少把背面也拉向正面外观的冲突。

优化完成后，辐射场可以从不同相机渲染；若要交付网格，还须提取表面并整理材质。多个正面、几何过平滑、细结构缺失或依赖视角的伪外观，都可能在环绕观察中出现。可渲染性没有证明表面可直接编辑，更没有证明生成背面等于真实玩具背面。

\subsection{单幅参考下的空间内容创建}
\label{sec:vision-multigen-single-reference}

单幅照片增加了可见轮廓、颜色和主体身份的约束，但背面和遮挡区域仍未知。创作前可约定正面保持原耳朵、眼睛与花纹，允许设计背面纹理。这样才能把偏离参考的错误，与合理补全但无法验证的内容区分开。

\subsubsection{Wonder3D的多视角颜色、法线与表面融合}
\label{sec:vision-multigen-wonder3d}

只生成多个方向的颜色图，重建时仍可能把纹理变化解释成凹凸。Wonder3D同时生成多视角颜色图和法线图，让外观与表面朝向互相约束。法线图的每个有效像素给出该处表面的朝向向量，它不直接给出深度或三维点位置\scite{vision-wonder3d2024}

模型以单幅参考图、目标相机和输出域条件为输入。输出域指当前生成的是颜色还是法线，由额外条件编码接入去噪网络。跨视角注意力让不同相机下的表示读取彼此信息，跨域注意力又让颜色与法线交换信息。仅有域标记能告诉网络输出哪类数据，却不能单独保证两类结果描述同一个表面。

训练使用三维资产渲染得到的对应图像与法线，微调具有图像条件能力的二维扩散模型。目标相机与输入参考之间的关系作为条件保留，采样时从多视角、多域噪声共同生成结果。原方法采用六个视角，未观测背面由生成先验补充。六幅看起来各自合理的图，仍可能在细耳朵、尾巴和遮挡位置互相矛盾。

形成几何时，方法优化一个\term{有符号距离场}（signed distance field，SDF）：空间点到表面的距离带有内外符号，零值集合定义表面。若约定物体外部为正，在可微且梯度非零处，
$\bm n(\bm p)=\nabla f(\bm p)/\|\nabla f(\bm p)\|_2$
给出向外单位法线。法线由SDF对空间的一阶导数得到；优化涉及法线的损失时，还要对场参数继续求导，不能把“求法线”本身称为计算SDF的二阶空间导数。

从各生成视图抽取像素射线，用SDF渲染颜色、法线和前景占据，再与相应生成图比较。颜色、法线与前景掩码分别提供外观、朝向和轮廓约束；梯度范数接近1的正则项保持距离场性质，稀疏与平滑项限制浮游表面和局部噪声。这些是对生成多视图的融合，不是把六张图直接粘成六面体。

生成法线也可能错误，因此原方法依据观察方向赋权。设$\bm v_i$是从相机射入表面的单位方向，$\bm n_i$是生成的向外单位法线，$\widehat{\bm n}_i$是当前渲染法线。几何法线损失为
\begin{equation}
 L_{\mathrm n}
 =\frac{\sum_i w_i(1-\widehat{\bm n}_i^\top\bm n_i)}
         {\sum_i w_i},
 \quad
 w_i=
 \begin{cases}
  0,&\bm v_i^\top\bm n_i>\eta,\\
  \exp(|\bm v_i^\top\bm n_i|),&\bm v_i^\top\bm n_i\leq\eta,
 \end{cases}
 \label{eq:vision-wonder-normal}
\end{equation}
其中$\eta$为接近0的负阈值，且当前批次要求$\sum_iw_i>0$。相机看到的向外表面法线应大致逆着入射视线，点积应为负；接近掠射的法线获得较低权重，朝向明显矛盾的结果被排除。全被排除时不能除以零，应跳过该项或重新取样。

取两个教学射线，方向点积为$-1$和$-0.5$，法线差异$1-\widehat{\bm n}_i^\top\bm n_i$分别为0.1和0.3，阈值取$-0.01$。两项权重为$e$和$e^{0.5}$，归一化损失约0.17551，比等权平均0.2更靠近正视方向的误差。它说明赋权的含义，不证明正视法线一定正确。颜色和掩码误差还采用剔除一部分最大残差的策略，降低局部不一致的影响，但可能同时忽略真实细节，须观察最终表面。

输入玩具狗正面时，可以检查输入侧轮廓是否保持，再环绕检查耳朵是否连接、纹理是否有接缝。若没有真实尺度或额外视图，背面纹理与绝对尺寸都不具有测量依据。用于展示的成果也许已经足够，用于动画绑定或制造则还要检查拓扑、厚度与材质等要求。

\subsection{多幅与连续观测下的空间内容形成}
\label{sec:vision-multigen-observed-space}

多幅观测能限制更多表面，前提是相机和时间关系正确。静态场景拟合及新视角合成已在\ref{sec:vision-generation-novel-view}节展开。若玩具狗正在运动，就不能把不同时间的腿部位置全部拟合为同一静态表面，需要一个随时间变化的场景表示。

\subsubsection{4D Gaussian Splatting的时空平面与高斯形变}
\label{sec:vision-multigen-4dgs}

Wu等人的4D Gaussian Splatting保留一组基准三维高斯，用位置和时间预测高斯形变，再渲染相应时刻。这里的“四维”指三维空间随时间变化，采用的是该文的时空平面与形变路线，不与其他同名的四维高斯参数化混用\scite{vision-4dgs2024}

静态高斯的中心、旋转、尺度、不透明度和颜色见\ref{sec:vision-generation-gaussians}节。形变模块读取基准中心$\bm\mu_i=(x_i,y_i,z_i)^\top$及查询时间$\tau$，在$(x,y)$、$(x,z)$、$(y,z)$、$(x,\tau)$、$(y,\tau)$和$(z,\tau)$六种二维平面上取特征。每个平面用四个邻近网格点双线性插值，因此不用存储完整四维稠密网格。

在同一分辨率下，六个平面特征逐元素相乘；不同分辨率的结果拼接，再经小型网络输出位置、旋转和尺度的变化。设$\mathcal P$为上述六组坐标对，$r$标记分辨率，可写
\begin{equation}
 \begin{aligned}
 \bm h_i(\tau)
 &=\mathop{\operatorname{concat}}_r
   \left[
    \mathop{\bigodot}_{(a,b)\in\mathcal P}
       \operatorname{interp}
        (\bm F_{a,b,r},a_i,b_i)
   \right],\\
 \Delta\bm\mu_i(\tau)&=D_{\mathrm{pos}}(\bm h_i(\tau))\text{。}
 \end{aligned}
 \label{eq:vision-4d-planes}
\end{equation}
当坐标对含$\tau$时，式中的相应坐标取查询时间；$\bm F_{a,b,r}$表示带特征通道的网格数组。相邻位置会共享网格参数，使附近高斯的形变具有共同表示，但不会强制它们属于同一个刚体。

中心按$\bm\mu_i(\tau)=\bm\mu_i+\Delta\bm\mu_i(\tau)$更新，旋转和尺度也在相应参数化中加入预测变化。渲染前还须把旋转参数恢复为合法旋转、尺度恢复为正值，不能把任意残差直接当作正定协方差。原路线的高斯不透明度与颜色参数仍来自基准集合，查询相机时再进行投影与透明度合成。

教学上，一个中心原为$(0,0,2)^\top$，在两个时刻预测的位移分别为$(0,0,0)^\top$和$(0.2,0,0)^\top$。若固定针孔相机焦距为100像素、主点横坐标为0，投影横坐标从0变为$100\times0.2/2=10$像素。纵深、旋转和尺度变化还会影响投影足迹，不能只把二维图中的点平移而忽略高斯形状。

训练先优化静态高斯获得初始场景，再联合学习基准表示与形变场。每次选择观测相机及时间，渲染当时的高斯，以对应图像作为监督，并对网格施加平滑正则。使用时输入任意查询相机和时间，生成该位置的视图；参数、平面特征和高斯集合共同构成需要保存的场景。

验收应把已观测时间、未观测视角和观测间时间分开。训练相机下拟合良好，可能仍在背面或快速运动时失败。单相机连续观测还把视角与时间交织在一起，某些表面变化难以唯一解释。高斯有基准编号也不等于它们具有真实物体的持久身份，形变网络更不是含质量、力和碰撞约束的物理模拟器。

\section{视觉条件下的文字与语音响应}
\label{sec:vision-multigen-response}

当用户询问“哪个物体先碰到桌面”，要生成的是基于视觉证据的回答及说出回答的声音。答案是否有依据属于\ref{chap:vision-multimodal-understanding}章的问题；本节关注文字内容如何转成语音，以及用户何时能听到它。这与给视频中的物体配音效具有不同目标。

\subsection{Qwen2.5-Omni的文字回答与流式语音形成}
\label{sec:vision-multigen-omni-response}

Qwen2.5-Omni的Thinker读取视觉、声音和问题，生成文字及其隐藏表示；Talker利用这些信息自回归地产生语音标记。结构与训练见\ref{sec:vision-models-qwen-omni}节。Talker同时接收文字词元信息，帮助区分语义接近但发音不同的词，而不只依据模糊的语义向量猜测音节。

语音标记需要解码才能播放。原方法把相邻标记分成块，以流匹配模型生成对应梅尔频谱，再由声码器恢复波形。解码当前块时读取有限的前后上下文，报告中的设置包含前面两块、当前块和后面一块。因此“流式”并不意味着拿到第一个文字词元就能零等待地发声；前看一块本身就需要一定积累。

响应延迟至少要区分首次发声和整段完成。设用户输入就绪时刻为0，第$k$块可播放音频准备好于$a_k$，播放时长为$d_k$，实际开始播放的时刻为$b_k$。忽略额外播放器缓冲时，
\begin{equation}
  b_1=a_1,\qquad
  b_k=\max(a_k,b_{k-1}+d_{k-1})\text{，}\quad k>1\text{。}
  \label{eq:vision-stream-playback}
\end{equation}
若音频尚未准备好，就必须等到$a_k$；若准备得很早，也要等前一块播完。这个排程关系与模型的具体参数无关，可以用来检查推理和播放是否衔接。

例如三块音频分别在0.8、1.2、1.7秒就绪，每块长0.5秒，则开始播放时刻为0.8、1.3、1.8秒，整段在2.3秒结束，中间无空隙。若第二块直到1.5秒才就绪，开始时刻变成0.8、1.5、2.0秒，出现0.2秒停顿。全部生成后播放在前一设置中要等到1.7秒才开始，并于3.2秒结束。以上均为教学排程，不是Qwen2.5-Omni的性能实测。

视觉问答还需决定是否可以提前说出某个结论。若视频中第二次接触尚未观察到，提前回答“花瓶先落桌”可能建立在猜测上。已经播放的词句无法像屏幕文字那样悄然撤回，纠正需要明确说出来。流式系统应在证据已经足够时输出相应内容，并把后来获得的信息与已说出的内容保持一致。

验收时同时查看视觉证据、显示文字和语音转写。漏读否定词、把“左边”说成“右边”、停顿切断关键数字，都可能使语音与正确文字不一致。自然度、可懂度、首次发声时间和最长停顿分别记录；更早开口不能补偿答错。该案例生成文字与语音回答，不承担视频内容生成。

\section{成果验收与研究方向}
\label{sec:vision-multigen-evaluation}

媒体成果既要独立可用，也要满足请求中规定的共同关系。自动特征分数、人类观看试听和具体任务检查承担不同职责。把它们保留为分项结果，才能判断失败发生在单项质量、条件执行还是输出之间的对应。

\subsection{单项成果的质量与可用性}
\label{sec:vision-multigen-quality}

视频可检查模糊、闪烁、运动和身份变化；声音可检查失真、声源、语言内容和可懂度；空间资产可检查表面、纹理、拓扑和编辑用途。指标应作用于相应交付对象。图像质量分数不包含声音，可旋转展示的辐射场也不自动满足网格编辑要求。

\subsubsection{VBench的分维度视频评价}
\label{sec:vision-multigen-vbench}

VBench把视频生成评价拆成主体与背景一致性、闪烁、运动平滑度、动态程度、画面质量及条件语义等维度，再为不同维度选择分析工具。其用途是定位差异，避免一个总分遮住主要失败\scite{vision-vbench2024}

主体一致性使用跨帧DINO特征相似度，背景一致性使用CLIP特征，动态程度借助RAFT光流估计。运动平滑度利用插帧模型的运动先验，条件语义则调用对象检测、动作分析或图文匹配等工具。这些工具输出的是代理测量，例如特征相似度高不能证明每根手指都保持正确。

评价流程从固定提示集合开始，为每条提示保留规定数量的生成样本，用同一版本的分项脚本处理，再与相应的人类判断比较。VBench用人类偏好标注检查自动测量的对应关系，但某个基准中的对应不保证在所有新场景中成立。工具对生成缺陷的误判也应保留为评价误差。

一个几乎不动的玩具狗片段可能在外观一致性上很好，却没有完成“绕花瓶走一圈”。另一个片段运动丰富，却在转身后变成另一只狗。二者需要分别报告动作完成与身份保持，不能把低动态当成稳定性的无条件优点。模型间比较还须固定帧率、长度、样本数和生成预算，否则评分差异可能来自不同制作条件。

\Needspace{65mm}
\subsubsection{VBench-2.0的物理与常识核验}
\label{sec:vision-multigen-vbench2}

更细的评价还要判断事件后果。VBench-2.0的2025年v2按人体表现、可控性、创造性、物理与常识组织测试，组合通用视觉语言模型、问题判定及专项异常检测器。它把“画面自然”进一步分解为可以检查的行为与状态\scite{vision-vbench22025}

在物理状态变化中，评价先根据提示形成具体的可见行为描述，再用多个问题核查生成视频。部分流程先排除初始状态不符合要求的片段，使状态变化判断对应原题。解释这种分数时，应同时报告初始条件的通过情况；只看过滤后的得分，会遗漏模型连题目起点都没有生成正确的失败。

对“玩具狗推倒花瓶”，教学核验可以分别问是否接触、花瓶是否倾斜、倾倒后是否保留同一个物体。仅有抬腿动作或一声碰撞都不足以回答这些问题。常识评价中的异常合并、突然复制与消失也需要跨帧观察，逐帧都像真图不表示帧间关系合理。

该基准的多视角一致性使用特征匹配等代理证据，因为生成视频通常没有对应真实三维模型。它不能直接测出真实几何误差。测试分数提高支持的是相应题目和评价流程下的改善，不是已经具备普遍物理模拟能力；自动判定仍须与可见事实及人类检查对照。

\subsection{条件与跨输出关系的满足}
\label{sec:vision-multigen-relations}

验收先把请求中的对象与关系落到同一份记录上。参考玩具狗的身份、杯子的材质、碰撞次数、声源及起声时刻相互关联，但应分别有可检查对象。表\ref{tab:vision-media-acceptance}列出贯穿示例中的核验方式，表中任务阈值需在生成前确定。

\begin{table*}[htbp]
 \centering
 \small
 \begin{tabularx}{\linewidth}{@{}p{25mm}p{44mm}X@{}}
  \toprule
  要求 & 主要证据 & 容易混淆的成功与失败 \\
  \midrule
  主体保持 & 原参考、转身前后与遮挡后画面 & 颜色相近仍可能改变耳朵、形状或身份 \\
  局部修改 & 原片段、编辑掩码、最终像素 & 指定区域改对，掩码外却有背景漂移 \\
  碰撞配声 & 接触时间、起声时间、事件次数 & 声音自然，但过早、过晚或多出一次 \\
  联合音画 & 共同事件、对白转写与口型 & 两路各自自然，却来自不同动作或声源 \\
  空间内容 & 参考覆盖、相机、表面与材质 & 新视图逼真，背面却无观测依据或不可编辑 \\
  语音回答 & 视觉依据、文字、语音与时序 & 文字正确，语音漏词，或早开口后长时间停顿 \\
  \bottomrule
 \end{tabularx}
 \caption{媒体制作的分项验收示例。证据应对应最终交付格式，掩码、相机、时间和事件来源需保存。各行回答不同问题，不能由某一项高分推断整件成果通过。}
 \label{tab:vision-media-acceptance}
\end{table*}
\FloatBarrier

跨输出关系可能随制作步骤改变。视频编辑把花瓶换成金属杯后，原陶瓷碰撞声可能已经不合适；三维资产改短了尾巴，沿用旧轮廓生成的多视角宣传图可能不再一致。应从受影响的关系出发重新核验，而不是只检查被修改的文件。

空间成果需标明哪些表面有输入观测支持，哪些由先验补出。可以在留出的真实视图上测量渲染误差；没有真实背面时，只能检查创作合理性、内部一致性及用途。把全部生成视角当作“验证视角”，会让同一个生成先验同时提供答案和检验依据。

\subsection{制作预算与任务完成程度}
\label{sec:vision-multigen-budget}

一条成功视频可能来自大量候选与多轮局部重做。比较制作能力时，应固定输入、时长、分辨率、模型版本、候选数量和编辑预算，并记录达到可用成果的总代价。只展示最佳样本会把筛选成本和失败概率都隐藏起来。

假设两个教学流程都交付一个四秒片段。甲生成4个候选，每个占用2分钟计算，再做一次1分钟的局部修改，累计计算量为9设备分钟；乙生成1个候选，占用5设备分钟。两者都通过当前验收，也不能因此说甲更省资源。若甲的4个候选并行生成，其等待时间可能短于9分钟，但设备累计时间仍是9分钟，人工观看与选择时间还要另计。

任务完成率的分母应是事先选定的请求，而不是最后留下的候选。可以记录规定预算内是否至少有一个结果通过，并同时给出每请求候选数、重做次数与失败类型。这样既描述了最终可用性，也保留模型一次生成的不稳定性。

这些失败引出不同研究方向。长时身份与场景漂移需要持续保存状态，局部修改需要协调新旧内容的边界，联合音画需要表达共同事件，多视角创建需要协调先验与几何，流式回答则要协调证据到达和语言播放。扩大主干或增加条件只提供更多表示和学习机会，究竟解决了哪个问题，仍要通过相应时间范围、输入条件和制作预算下的结果回答。

\section{本章小结}
\label{sec:vision-multigen-summary}

媒体制作的起点是规定什么已知、什么允许改变，以及哪些关系必须保持。整段时空处理与潜空间生成帮助形成视频，掩码和原片段限制编辑范围，已知时间位置约束续接与补全。固定视频配声把画面作为条件，联合生成则让声音和画面共同变化；二者都要把最终事件放回同一时间轴核验。

空间生成进一步要求不同视角共同解释同一个对象。二维先验可以通过渲染梯度或多视角颜色与法线提供约束，多幅观测则让场景拟合具有更多证据。无论表示多么自然，观测支持的表面与先验补出的内容仍须区分。视觉条件响应还把生成延伸到播放过程，文字正确、语音可懂和及时作答共同决定使用效果。

开篇的玩具狗片段只有在身份、动作、接触、声音和时间关系共同满足时才算完成。评价因而既看单项质量，也看条件与跨输出关系，并把候选筛选和修改计入预算。可用成果来自这些约束的共同落实，不能由一幅漂亮画面、一条自然声轨或一个汇总分替代。
